\documentclass[twoside]{projektInzynierskiMS1z}
\usepackage{polski}
\usepackage[utf8]{inputenc}
\usepackage{amsmath}
\usepackage{graphicx}
\usepackage{url}
\usepackage{setspace}
\usepackage{tikz}
\usepackage{tikz}
\usepackage{dirtree}
\usepackage[colorlinks=true,linkcolor=black,citecolor=black,urlcolor=black]{hyperref}
\usetikzlibrary{shapes.geometric, arrows.meta, positioning, fit, backgrounds, calc}
\usetikzlibrary{arrows.meta, shapes.geometric}
%\drukJednostronny

%% naglowek bibliografii po polsku
\renewcommand{\refname}{Literatura}

%% tytul, promotor i autorzy
\title{Projekt i implementacja systemu symulacji oraz automatycznego
zarządzania ruchem kolejowym w środowisku webowym}

\promotor{dr inż. Rafał Brociek}
\rodzajProm{PROWADZĄCY PRACĘ}

\specjalnosc{Inżynieria Analizy Danych}
\ustawKierunekW{informatyka}
\jednostkaUczelni{WYDZIAŁ MATEMATYKI STOSOWANEJ}
\miasto{GLIWICE}
\rok{2026}

%% kazdy autor: imie nazwisko, nr albumu, opis wkladu
\autor{Martyna SZANDAŁA}{311093}{Opis wkładu}
\autor{Natalia TOMALA}{311109}{Opis wkładu}
\autor{Piotr DUSIŃSKI}{311489}{Opis wkładu}

%%%%%%%% zestaw komend opcjonalnych
%\NumeryNaPoczatku
%\subsectionWzory
%\sectionWzory
%\rozdzialy
%\literowaNumeracjaDodatkow
%\rzymskaNumeracjaDodatkow

\newtheorem{tw}{Twierdzenie}
\newtheorem{twa}{Twierdzenie}
\newtheorem{dd}{Definicja}

%%%%%%%%%%%%%%%%%%%%%%%%%%% strona druga -- streszczenia
\drugaStrona

\strDrNap{Streszczenie}

Przedmiotem niniejszej pracy jest zaprojektowanie oraz implementacja
wielowarstwowej aplikacji webowej służącej do interaktywnej symulacji
i monitorowania wirtualnego ruchu kolejowego. Głównym celem projektu było stworzenie narzędzia, które w przystępny sposób obrazuje złożone procesy zarządzania ruchem kolejowym, w tym wyznaczanie tras pociągów, reagowanie na awarie infrastruktury (zamknięcia odcinków, ograniczenia prędkości) oraz automatyczne wyznaczanie objazdów.
\strDrNap{Słowa kluczowe}

symulacja ruchu kolejowego, aplikacja webowa, baza grafowa, algorytm A*,
Memgraph, rerouting

\strDrNap{Thesis title}

Design and implementation of a web-based railway traffic simulation
and automatic management system

\strDrNap{Abstract}

The subject of this thesis is the design and implementation of a multi-layered
web application for the interactive simulation and monitoring of virtual railway
traffic. The main goal of the project was to create a tool that illustrates, in an
accessible way, the complex processes of railway traffic management, including
train routing, responses to infrastructure failures (track section closures and
speed restrictions) and automatic rerouting.

\strDrNap{Keywords}

railway traffic simulation, web application, graph database, A* algorithm,
Memgraph, rerouting

\koniecDrugaStrona

\begin{document}

%% interlinia 1,5 dla tekstu podstawowego (zgodnie z formularzem Z4-PU12)
\onehalfspacing

%% UWAGA: naglowek "Wstep" (rozdzial 1) jest generowany automatycznie przez
%% klase zaraz po spisie tresci -- dlatego tu NIE wpisujemy \section{Wstep}.
%% Ponizszy tekst stanowi tresc wstepu.

Niniejszy projekt inżynierski koncentruje się na opracowaniu i implementacji
interaktywnej aplikacji webowej służącej do symulacji oraz monitorowania
wirtualnego ruchu kolejowego. Realizacja projektu obejmowała przygotowanie
koncepcji logiki biznesowej, zaprojektowanie interfejsu graficznego oraz
wdrożenie wielowarstwowej architektury systemu. Głównym celem aplikacji jest
stworzenie środowiska (przypominającego interaktywną makietę kolejową), które
w sposób przystępny i wizualny obrazuje skomplikowane procesy zarządzania ruchem
pociągów, takie jak wyznaczanie tras pociągów czy reagowanie na awarie infrastruktury.

W pracy szczegółowo opisano cały proces realizacji projektu -- od analizy
wymagań funkcjonalnych, przez projektowanie architektury opartej na strukturach
grafowych i interfejsu, po implementację wybranych algorytmów nawigacyjnych
i testowanie systemu. Ważnym elementem cyklu życia projektu było zastosowanie
zwinnych metodyk wytwarzania oprogramowania (ang. \emph{Agile}), co umożliwiło
elastyczne planowanie i etapowe rozszerzanie skali symulacji. Zrealizowane
rozwiązanie jest skierowane zarówno do pasjonatów kolejnictwa, jak i studentów
kierunków związanych z transportem, stanowiąc narzędzie ułatwiające zrozumienie, jak awarie infrastruktury wpływają na ruch pociągów i w jaki sposób wyznaczane są trasy objazdowe. Projekt
ten stanowi również solidną podstawę do dalszej rozbudowy o zewnętrzne źródła
danych.

\subsection{Geneza}

W ostatnich latach można zaobserwować silny trend cyfryzacji w sektorze transportu
publicznego, w tym wdrażanie koncepcji ,,cyfrowych bliźniaków'' (ang. \emph{digital
twins}) do symulowania zachowań rzeczywistej infrastruktury~\cite{digitaltwin}.
Wraz z tym zjawiskiem rośnie zapotrzebowanie na narzędzia edukacyjne i analityczne,
które w czytelny sposób tłumaczą zasady bezpieczeństwa na kolei. Inspiracją do
stworzenia aplikacji była obserwacja, że procesy decyzyjne dyspozytorów oraz
mechanizmy zabezpieczające tory są dla przeciętnego pasażera całkowicie niewidoczne,
a ich zrozumienie z samych podręczników jest trudne. Istniejące na rynku
i ogólnodostępne rozwiązania posiadają w tym zakresie pewne luki.

Popularne serwisy dla pasażerów (np. Portal Pasażera PKP PLK) skupiają się wyłącznie
na prezentacji aktualnych opóźnień i rozkładów jazdy, nie dając wglądu w przyczyny
tych zjawisk ani nie posiadając warstwy symulacyjnej. Z kolei zaawansowane symulatory
kolejowe (takie jak MaSzyna) kładą nacisk na trójwymiarowy realizm i perspektywę
maszynisty prowadzącego pojedynczy pojazd. Wymagają one wydajnego sprzętu i nie
pozwalają na spojrzenie na sieć kolejową z perspektywy makro -- jako na złożony
system logistyczny. Dostępne amatorskie symulatory urządzeń sterowania ruchem
(np. ISDR) są natomiast programami komputerowymi o wysokim progu wejścia, opartymi
na skomplikowanych, technicznych interfejsach odzwierciedlających prawdziwe pulpity
kostkowe. Brakowało na rynku nowoczesnego, przeglądarkowego narzędzia, które
łączyłoby rzeczywiste lub symulowane dane z przystępną symulacją na mapie
dwuwymiarowej.

\subsection{Odniesienie do rynku}

Rynek systemów informatycznych dla kolei (ang. \emph{Railway Management Systems})
dynamicznie rośnie, co jest napędzane potrzebą modernizacji infrastruktury
i zwiększania przepustowości linii. Prognozy rynkowe wskazują, że globalna wartość
tego sektora będzie w najbliższych latach systematycznie rosnąć, wymuszając
kształcenie nowych kadr w branży logistycznej~\cite{railwaymarket}.

Obecnie na rynku profesjonalnym dominują zaawansowane, zamknięte systemy
wykorzystywane przez zarządców infrastruktury (np. System Wspomagania Dyżurnego
Ruchu -- SWDR). Są to jednak narzędzia komercyjne, do których dostęp mają wyłącznie
uprawnieni pracownicy kolei. Z drugiej strony, w sektorze edukacyjnym i akademickim
brakuje lekkich, darmowych rozwiązań pozwalających na testowanie scenariuszy
kryzysowych na torach.

Większość dostępnych projektów akademickich skupia się na suchej analizie danych
algorytmicznych, pomijając aspekt wizualny. Zrealizowana aplikacja wypełnia tę niszę.
   Zapewnia ona środowisko symulacyjne oparte na rzeczywistej sieci kolejowej
   węzła śląskiego (57 stacji i 71 odcinków), w którym każdy może obserwować ruch
   pociągów, a zalogowany użytkownik może celowo wywoływać awarie i badać reakcję
   systemu.

\subsection{Cechy wyróżniające aplikację}

Projekt łączy intuicyjny interfejs z zaawansowaną architekturą logiki biznesowej.
Do głównych cech wyróżniających aplikację na tle prostych wizualizacji należą:

\begin{itemize}
\item \textbf{Zastosowanie grafowej bazy danych:} Wykorzystanie bazy działającej
w pamięci (Memgraph) pozwoliło na wysoce wydajne odwzorowanie sieci kolejowej.
Zapisanie stacji jako węzłów, a torów jako krawędzi, umożliwiło naturalne
i błyskawiczne operacje na strukturze sieciowej, niezbędne przy symulacjach
o dużej skali.

\item \textbf{Dynamiczny rerouting z algorytmem A*:} Aplikacja posiada wbudowany
silnik decyzyjny. Gdy odcinek toru zostaje zamknięty, pociągi, których trasa
przez niego prowadzi, automatycznie otrzymują trasę alternatywną. Wyznacza ją
zaimplementowany od podstaw algorytm A*, który szuka trasy o najkrótszym czasie
przejazdu. Heurystyką jest czas pokonania odległości ortodromicznej (obliczonej
wzorem Haversine'a) z prędkością maksymalną pociągu. Oszacowanie to nigdy nie
przekracza rzeczywistego czasu przejazdu, co gwarantuje optymalność wyznaczonej
trasy.

\item \textbf{Niezależność od zewnętrznych źródeł danych:} W odróżnieniu od
systemów zależnych od zewnętrznych interfejsów API projekt zrealizowano jako
autonomiczną piaskownicę. Sieć kolejowa (57 stacji, 71 odcinków) i flota
52 pociągów zostały wygenerowane z rozkładu jazdy Kolei Śląskich w formacie
GTFS~\cite{gtfsks} oraz wykazów Regulaminu sieci PKP PLK~\cite{plkregulamin}
i zapisane jako dane startowe bazy. Zapewnia to powtarzalność testów,
uniezależnia symulację od dostępności zewnętrznych usług i pozwala etapowo
rozszerzać zbiór danych.

\item \textbf{Model ruchu i stanów infrastruktury:} System nie tylko przesuwa
znaczniki na mapie -- każdy pociąg porusza się zgodnie z prędkościami
maksymalnymi odcinków pochodzącymi z wykazów PKP PLK, a każdy odcinek może być
czynny, objęty ograniczeniem prędkości lub zamknięty. Zmiana stanu odcinka
wpływa na ruch pociągów: powoduje zwolnienie, zatrzymanie lub wyznaczenie
objazdu. Dzięki temu aplikacja stanowi stabilne środowisko do przeprowadzania
symulacji scenariuszy awaryjnych.
\end{itemize}

\subsection{Wykorzystanie w praktyce}

Stworzona aplikacja może znaleźć szerokie zastosowanie w wielu realnych scenariuszach
edukacyjnych, demonstracyjnych i analitycznych, do których należą między innymi:

\begin{itemize}
\item \textbf{Prowadzenie zajęć z logistyki i transportu:} wizualna demonstracja
w czasie rzeczywistym, jak zamknięcie odcinka lub ograniczenie prędkości zmienia
trasy i czasy przejazdu pociągów w sieci.

\item \textbf{Analiza zarządzania kryzysowego typu what-if:} Dzięki funkcji ręcznego zgłaszania awarii, użytkownik zyskuje możliwość obserwacji,
jak zamknięcie jednego, kluczowego odcinka na Śląsku wpływa na trasy i czasy przejazdu pociągów w całej sieci oraz jak reaguje algorytm wyznaczający objazdy.

\item \textbf{Porównanie kategorii pociągów:} pociągi dalekobieżne (160~km/h),
regionalne (120~km/h, z doliczanym czasem postojów) i towarowe (80~km/h) mają
różne parametry, dlatego algorytm A* może wyznaczyć dla nich różne trasy między
tymi samymi stacjami. Pozwala to pokazać, że najkrótsza trasa nie zawsze jest
najszybsza.

\item \textbf{Zrozumienie algorytmiki grafowej:} Aplikacja jest doskonałym narzędziem
dydaktycznym dla studentów informatyki, obrazującym praktyczne użycie bazy Memgraph
i algorytmu A* na rzeczywistych danych przestrzennych.
\end{itemize}

Dzięki swojej architekturze opartej na przejrzystym interfejsie graficznym aplikacja
obniża próg wejścia w złożony świat zarządzania ruchem kolejowym, pozwalając na obrazowe i
bezpieczne eksperymentowanie zarówno osobom niezwiązanym z koleją, jak i logistykom w uproszczony sposób.

\section{Założenia projektowe aplikacji}

Fazę projektową poprzedziła dogłębna analiza dostępnych rozwiązań
oraz zdefiniowanie celów systemowych w oparciu o zasady inżynierii
wymagań~\cite{sommerville2016}. Ponieważ tworzenie zaawansowanego symulatora wymaga
precyzyjnego zbalansowania wydajności obliczeniowej z realizmem odwzorowania procesów,
zespół projektowy przyjął zwinne podejście do cyklu wytwarzania oprogramowania
(ang. \emph{Agile}). Metodyka ta pozwala na iteracyjne dostarczanie kolejnych
modułów funkcjonalnych oraz elastyczne reagowanie na problemy techniczne pojawiające
się w trakcie tworzenia projektu~\cite{agilemanifesto}.

Prace projektowe podzielono na etapy, zaczynając od definicji logiki biznesowej
i wymagań krytycznych, a kończąc na planowaniu modułów opcjonalnych. Taki podział
jest zgodny z dobrymi praktykami projektowania systemów informatycznych, gdzie
priorytetyzacja zadań pozwala na zachowanie spójności architektury przy rosnącej
skali projektu~\cite{martin2017clean}.

\subsection{Klasyfikacja wymagań systemowych}

Proces specyfikacji pozwolił na jasne wyodrębnienie rdzenia systemu (ang. \emph{core})
od funkcji rozszerzających. Zgodnie z teorią projektowania systemów zorientowanych na
dane, wymagania podzielono na kategorie uwzględniające ich wpływ na stabilność silnika
symulacyjnego:

\begin{itemize}
\item \textbf{Wymagania krytyczne:} Funkcjonalności niezbędne do realizacji
podstawowego modelu ruchu. Należą do nich: interaktywna wizualizacja węzłów
i krawędzi sieci, mechanizm zasilania danymi, ruch pociągów po wyznaczonych
trasach oraz model stanów odcinków (odcinek czynny, z ograniczeniem prędkości
lub zamknięty). Bez tych elementów system nie mógłby pełnić roli symulatora
ruchu kolejowego.

\item \textbf{Wymagania rozszerzone:} Moduły podnoszące wartość analityczną projektu,
takie jak algorytm dynamicznego wyznaczania objazdów (rerouting) oraz moduł zdarzeń losowych wraz z możliwością ręcznego zgłaszania incydentów. Ich implementacja pozwala na symulowanie rzeczywistych problemów występujących w sieciach transportowych,
wzbogacając znacząco wartość dydaktyczną całego projektu~\cite{sommerville2016}.

\item \textbf{Wymagania opcjonalne:} Funkcje o niższym priorytecie, których realizacja
jest uzależniona od wydajności końcowej wersji systemu (np. integracja z zewnętrznymi
interfejsami czy import rozbudowanych rozkładów jazdy). Ograniczenia projektu obejmują
świadome wykluczenia (np. rezygnacja z renderowania mapy całego kraju), mające na celu
utrzymanie wysokiej responsywności aplikacji webowej wyłącznie w środowisku przeglądarkowym.
\end{itemize}

\subsection{Opis kluczowych funkcjonalności systemu}

Na podstawie przeprowadzonej analizy wyodrębniono moduły funkcjonalne stanowiące trzon
aplikacji. Zespół zdecydował się na architekturę opartą na niezależnym silniku
symulacyjnym, co wymagało zaprojektowania następujących mechanizmów:

\subsubsection{Silnik symulacyjny i stany odcinków}

Sercem systemu jest moduł, który w stałych odstępach czasu (domyślnie co 1~s)
przelicza pozycje pociągów. Zastosowano uproszczony model kinematyczny: pociąg
porusza się po odcinku ze stałą prędkością równą mniejszej z dwóch wartości --
prędkości maksymalnej odcinka i prędkości maksymalnej pociągu. Postęp pociągu na
odcinku o długości $L$ [km] w jednym kroku symulacji wynosi
\begin{equation}
p_{k+1} = \min\left(1,\; p_k + \frac{v \cdot \Delta t_{\mathrm{sym}}}{3600 \cdot L}\right),
\qquad v = \min\left(v^{\mathrm{odc}}_{\max},\, v^{\mathrm{poc}}_{\max}\right),
\end{equation}
gdzie $v$ wyrażono w km/h, a $\Delta t_{\mathrm{sym}}$ to czas symulowany
przypadający na jeden krok (iloczyn interwału kroku, skali czasu równej 60 oraz
wybranego przez użytkownika mnożnika tempa od $0{,}5$ do $2$).

Każda skierowana krawędź grafu ma jeden z trzech stanów: \texttt{active} (ruch
bez ograniczeń), \texttt{restricted} (ograniczenie prędkości do połowy wartości
maksymalnej, nie mniej niż 20~km/h) oraz \texttt{blocked} (odcinek zamknięty).
Stany te ustawiają zdarzenia -- losowe lub zgłoszone ręcznie. Pociąg nie wjeżdża
na odcinek zamknięty, a jeśli zamknięcie nastąpi w trakcie przejazdu, zostaje
zatrzymany. Mechanizm ten w uproszczeniu odpowiada zamknięciom torowym stosowanym
w eksploatacji kolei. Model nie obejmuje natomiast blokady zajętości odstępów,
która w urządzeniach sterowania ruchem kolejowym (SRK) gwarantuje, że w jednym
odstępie znajduje się tylko jeden pociąg~\cite{srk}; jej implementację
przewidziano jako kierunek dalszego rozwoju.

\subsubsection{Dynamiczny rerouting i algorytmy ścieżkowe}

Statyczne przypisanie tras okazało się niewystarczające w warunkach dynamicznych
awarii. Zastosowano więc mechanizm automatycznego wyznaczania tras alternatywnych
oparty na algorytmie A*. Jest to algorytm przeszukiwania heurystycznego, który
przy dopuszczalnej heurystyce gwarantuje znalezienie ścieżki o minimalnym koszcie,
rozwijając zwykle mniej wierzchołków niż algorytm Dijkstry~\cite{astar1968}.

Kosztem krawędzi $(u,w)$ jest czas przejazdu odcinka o długości $L_{uw}$
z prędkością $v=\min(v^{\mathrm{odc}}_{\max}, v^{\mathrm{poc}}_{\max})$,
powiększony dla pociągów regionalnych o 1~min za każdy postój pośredni
($k_{\mathrm{post}}$). Heurystyka szacuje czas dojazdu ze stacji $n$ do stacji
docelowej $t$:
\begin{equation}
c(u,w) = \frac{60\,L_{uw}}{v} + k_{\mathrm{post}}, \qquad
h(n) = \frac{60\, d_{\mathrm{hav}}(n, t)}{v^{\mathrm{poc}}_{\max}},
\end{equation}
gdzie $d_{\mathrm{hav}}$ to odległość ortodromiczna obliczona wzorem
Haversine'a~\cite{haversine}. Odcinki zamknięte są pomijane, a na odcinkach
z ograniczeniem prędkości przyjmowana jest prędkość obniżona.

Heurystyka $h(n)$ jest dopuszczalna, ponieważ długość odcinka mierzona wzdłuż
toru nie jest mniejsza od odległości ortodromicznej jego końców, a pociąg nie
może jechać szybciej niż $v^{\mathrm{poc}}_{\max}$. Pierwszy warunek
zweryfikowano dla wszystkich 71 odcinków sieci -- najmniejszy iloraz długości
torowej do ortodromicznej wynosi $1{,}001$. Z nierówności trójkąta dla odległości
ortodromicznej wynika ponadto spójność (monotoniczność) heurystyki.

\subsubsection{Zgłaszanie incydentów i analiza ,,what-if''}

Aby aplikacja miała charakter narzędzia badawczego, oprócz zdarzeń losowych,
generowanych automatycznie zgodnie z procesem Poissona, wprowadzono możliwość
ręcznego zgłaszania incydentów (awaria sieci trakcyjnej, wykolejenie,
ograniczenie prędkości, awaria sterowania ruchem) dla wskazanego odcinka,
stacji lub pociągu. Funkcja jest dostępna dla użytkowników z uprawnieniem
\texttt{simulation.control} -- domyślnie dla ról ,,użytkownik'' i
,,administrator'' -- natomiast gość ma wyłącznie podgląd symulacji. Pozwala to
na prowadzenie testów scenariuszowych (ang. \emph{what-if}), co jest standardową
praktyką w projektowaniu systemów krytycznych, pozwalającą na identyfikację tak
zwanych wąskich gardeł infrastruktury jeszcze przed ich wystąpieniem
w rzeczywistości~\cite{sommerville2016}.

\section{Instrukcja obsługi programu}

bla bla jadi jadi jada
% Pierwsze co widzi użyktownik
\subsection{Strona prezentacyjna} 
% Faktyczny panel symulacji
\subsection{Panel symulacyjno-projektowa}
% Panel do zarządzania użykownikami oraz rolami
\subsection{Panel administracyjna}

\section{Opis architektury programu}
W niniejszym rozdziale opisano architekturę aplikacji oraz zastosowane wzorce
projektowe. Architektura systemu opiera się na trzech głównych warstwach:
\begin{itemize}
    \item Warstwa prezentacyjna (ang. \emph{frontend}) -- interfejs użytkownika
    wraz ze wszystkimi elementami interaktywnymi aplikacji. Za jego
    pośrednictwem użytkownik wykonuje operacje przewidziane w funkcjonalności
    programu.
    \item Warstwa komunikacyjno-zarządzająca (ang. \emph{backend}) -- warstwa
    działająca po stronie serwera, niewidoczna dla użytkownika końcowego.
    Odpowiada za logikę aplikacji, komunikację z pozostałymi warstwami,
    obliczenia oraz uwierzytelnianie i autoryzację.
    \item Warstwa bazodanowa -- służy do trwałego przechowywania danych
    aplikacji, m.in. kont użytkowników i ról, struktury sieci kolejowej (stacji
    i odcinków), stanu pociągów, zdarzeń oraz scenariuszy rozkładów.
\end{itemize}
Fundamentem tych warstw jest środowisko kontenerowe, które zapewnia ich izolację
i komunikację sieciową oraz znacząco ułatwia pracę zespołową. W kolejnych
podrozdziałach szczegółowo opisano komponenty składające się na całość
oprogramowania.

% Opisanie dockera i jego sieci-podsieci
\subsection{Środowisko}
Powszechnym standardem we współczesnej inżynierii oprogramowania jest dekompozycja aplikacji na niezależne komponenty, uruchamiane w postaci izolowanych kontenerów~\cite{kontener.microsoft}. W niniejszym podrozdziale szczegółowo omówiono architekturę oraz logikę środowiska uruchomieniowego projektu.

\subsubsection{Architektura stosu kontenerowego i segmentacja sieciowa}
% [X] 1. Usługi: memgraph-db, memgraph-lab, backend, frontend.
% [X] 2. Dwie sieci wirtualne (frontend-net, backend-net) i rolę backendu (dual-homed).
% [X] 3. Wewnętrzny DNS Dockera i Service Discovery (komunikacja po nazwach usług).
% [X] 4. Trwałość danych (Data Persistence) – nazwany wolumen db-data dla bazy grafowej.                   % [X] 5. [RYSUNEK / SCHEMAT ARCHITEKTURY KONTENEROWEJ]   

Architektura została oparta na narzędziu konteneryzacyjnym Docker. Stos usług
(ang. \emph{Compose stack}) zdefiniowany w pliku \texttt{compose.yaml} obejmuje
cztery usługi połączone dwiema sieciami wirtualnymi, które izolują od siebie
poszczególne warstwy. Kontener \texttt{frontend}, działający wyłącznie w sieci
\texttt{frontend-net}, nie ma bezpośredniego dostępu do bazy danych.
Pośrednikiem między interfejsem a bazą jest \texttt{backend}, podłączony
jednocześnie do sieci \texttt{frontend-net} i \texttt{backend-net}. Izolacja ta
dotyczy komunikacji między kontenerami: w obecnej, deweloperskiej konfiguracji
porty bazy danych (\texttt{7687}, \texttt{7444}) oraz Memgraph Lab
(\texttt{3000}) są dodatkowo opublikowane na maszynie hosta, co ułatwia pracę
programistom. W środowisku produkcyjnym porty te należałoby pozostawić
niedostępne z zewnątrz. Komunikacja wewnątrz sieci nie wymaga definiowania
statycznych adresów IP -- środowisko wykorzystuje wbudowany serwer DNS Dockera,
dzięki czemu usługi odwołują się do siebie po nazwach (np.
\texttt{backend:8000}). Stos obejmuje następujące usługi:
\begin{itemize}
    \item \texttt{memgraph-db} -- baza danych Memgraph uruchamiana z obrazu
    \texttt{memgraph/memgraph-mage} (Memgraph z biblioteką algorytmów MAGE),
    działająca w sieci \texttt{backend-net}. Nasłuchuje na portach
    \texttt{7687} (protokół Bolt) oraz \texttt{7444} (strumień logów dla
    Memgraph Lab). Nazwę i hasło użytkownika bazy pobiera ze zmiennych
    zdefiniowanych w pliku \texttt{.env}. Korzysta z nazwanego woluminu
    \texttt{db-data} zamontowanego w katalogu \texttt{/var/lib/memgraph}, co
    zapewnia trwałość danych (ang. \emph{data persistence}) i chroni stan grafu
    przed utratą po usunięciu kontenera.
    \item \texttt{memgraph-lab} -- graficzne środowisko Memgraph Lab wspierające
    pracę deweloperską: umożliwia wykonywanie zapytań i prezentuje wyniki,
    o ile to możliwe, w postaci grafu węzłów i relacji. Jest dostępne przez port
    \texttt{3000}, działa w sieci \texttt{backend-net} i uruchamiane jest po
    usłudze \texttt{memgraph-db}, z którą łączy się automatycznie.
    \item \texttt{backend} -- serwer aplikacji FastAPI podłączony do sieci
    \texttt{backend-net} i \texttt{frontend-net}, dostępny na porcie
    \texttt{8000}. Z pliku \texttt{.env} wczytuje dane dostępowe do bazy oraz
    parametry konta administratora. Uruchamiany jest po usłudze
    \texttt{memgraph-db}.
    \item \texttt{frontend} -- aplikacja SvelteKit działająca w sieci
    \texttt{frontend-net}, dostępna na porcie \texttt{5173} i uruchamiana po
    usłudze \texttt{backend}. Przeglądarka łączy się z backendem (również przez
    WebSocket) przez port \texttt{8000} opublikowany na hoście, natomiast
    renderowanie po stronie serwera (SSR) korzysta z adresu
    \texttt{backend:8000} wewnątrz sieci Dockera.
\end{itemize}

Kompletną architekturę stosu kontenerowego wraz z podziałem na sieci i przepływem danych zilustrowano na rysunku~\ref{fig:docker_architektura_full}.

\begin{figure}[p]
\centering
\resizebox{0.92\textwidth}{!}{%
\begin{tikzpicture}[
    font=\sffamily,
    >=Stealth,
    % Styl stacji roboczej hosta
    hostbox/.style={
        draw=gray!70!black, fill=gray!8, rounded corners=6pt, line width=1.1pt,
        minimum width=13.5cm, minimum height=1.3cm, align=center
    },
    % Styl standardowego kontenera
    appbox/.style={
        draw=blue!70!black, fill=white, rounded corners=5pt, line width=1.1pt,
        minimum width=3.8cm, minimum height=1.7cm, align=center, inner sep=6pt
    },
    % Styl backendu (most sieciowy)
    bridgebox/.style={
        draw=green!60!black, fill=green!5, rounded corners=6pt, line width=1.3pt,
        minimum width=5.4cm, minimum height=1.8cm, align=center, inner sep=6pt
    },
    % Styl bazy danych
    dbbox/.style={
        cylinder, draw=orange!80!black, fill=orange!8, shape border rotate=90, aspect=0.25,
        minimum width=3.6cm, minimum height=1.8cm, line width=1.2pt, align=center, inner sep=4pt
    },
    % Styl wolumenu
    volbox/.style={
        cylinder, draw=gray!70!black, fill=gray!15, shape border rotate=90, aspect=0.3,
        minimum width=2.0cm, minimum height=1.4cm, line width=1pt, align=center, font=\footnotesize
    },
    % Etykiety na połączeniach (zabezpieczone białym tłem)
    edgelbl/.style={
        font=\scriptsize, fill=white, draw=gray!40, rounded corners=3pt, inner sep=2.5pt, align=center
    },
    % Linie połączeń
    flow/.style={
        ->, line width=1.1pt, draw=gray!70!black
    },
    flowbolt/.style={
        <->, line width=1.3pt, draw=blue!70!black
    }
]

    % -------------------------------------------------------------
    % 1. RAMKA DOCKERA I SIECI (TŁO)
    % -------------------------------------------------------------
    % Główny kontener Docker Compose
    \draw[fill=gray!2, draw=gray!40, rounded corners=12pt, line width=0.9pt] 
        (-7.0, -2.0) rectangle (7.0, -15.5);
    % Tytuł w lewym dolnym rogu (całkowicie wolna przestrzeń!)
    \node[anchor=south west, font=\footnotesize\bfseries, text=gray!60!black] 
        at (-6.7, -15.3) {Docker Compose Stack};

    % Sieć: frontend-net (lewa strona)
    \draw[fill=cyan!4, draw=cyan!70!black, dashed, rounded corners=8pt, line width=1pt] 
        (-6.6, -2.7) rectangle (-0.3, -9.8);
    \node[anchor=north west, font=\footnotesize\bfseries, text=cyan!70!black] 
        at (-6.3, -2.9) {\texttt{frontend-net}};

    % Sieć: backend-net (prawa strona)
    \draw[fill=purple!4, draw=purple!70!black, dashed, rounded corners=8pt, line width=1pt] 
        (0.3, -2.7) rectangle (6.6, -15.0);
    \node[anchor=north east, font=\footnotesize\bfseries, text=purple!70!black] 
        at (6.3, -2.9) {\texttt{backend-net}};

    % -------------------------------------------------------------
    % 2. BLOK HOSTA (U GÓRY)
    % -------------------------------------------------------------
    \node[hostbox] (host) at (0, 0) {
        \textbf{\large Stacja robocza hosta / Przeglądarka internetowa} \\[1mm]
        \footnotesize Połączenia z zewnątrz kierowane do opublikowanych portów maszyny (\texttt{localhost})
    };

    % -------------------------------------------------------------
    % 3. KONTENERY
    % -------------------------------------------------------------
    % Frontend (w lewej kolumnie)
    \node[appbox, draw=cyan!80!black] (fe) at (-3.8, -4.8) {
        \textbf{\large frontend} \\[1mm]
        \small SvelteKit / Node.js\\
        \footnotesize Port: \texttt{5173}
    };

    % Memgraph Lab (w prawej kolumnie)
    \node[appbox, draw=purple!80!black] (lab) at (3.8, -4.8) {
        \textbf{\large memgraph-lab} \\[1mm]
        \small Panel GUI bazy\\
        \footnotesize Port: \texttt{3000}
    };

    % Backend (centralny most na granicy sieci)
    \node[bridgebox] (be) at (0, -8.2) {
        \textbf{\large backend (FastAPI)} \\[1mm]
        \small Brama / Most obu sieci\\
        \footnotesize Port kontenera: \texttt{8000}
    };

    % Baza danych Memgraph
    \node[dbbox] (mg) at (2.0, -12.8) {
        \textbf{\large memgraph-db} \\[1mm]
        \small Baza grafowa Memgraph\\
        \footnotesize Porty: \texttt{7687}, \texttt{7444}
    };

    % Wolumen (z bezpiecznym odstępem)
    \node[volbox] (vol) at (5.5, -12.8) {
        \textbf{db-data} \\[1mm]
        \scriptsize Wolumen
    };

    % -------------------------------------------------------------
    % 4. STRZAŁKI POŁĄCZEŃ (CZYSTE KORYTARZE)
    % -------------------------------------------------------------
    % Host -> Frontend (etykieta powyżej przerywanej linii sieci!)
    \draw[flow] (-3.8, -0.65) -- (fe.north)
        node[pos=0.32, edgelbl] {HTTP :5173};

    % Host -> Backend (centralny korytarz)
    \draw[flow] (0, -0.65) -- (be.north)
        node[pos=0.42, edgelbl] {REST / WebSocket :8000};

    % Host -> Memgraph Lab (etykieta powyżej przerywanej linii sieci!)
    \draw[flow, dashed] (3.8, -0.65) -- (lab.north)
        node[pos=0.32, edgelbl] {GUI :3000};

    % Frontend -> Backend (SSR)
    \draw[flow] (fe.south) |- (-2.7, -8.2)
        node[pos=0.25, left=2pt, edgelbl] {SSR (:8000)};

    % Backend -> Memgraph DB (Poprowadzona środkiem, etykieta w pustej przestrzeni!)
    \draw[flowbolt] (be.south) -- ++(0, -1.6) -| ([xshift=-10pt]mg.north)
        node[pos=0.35, above=2pt, edgelbl] {Bolt Cypher :7687};

    % Memgraph Lab -> Memgraph DB (Wchodzi czysto z prawej strony)
    \draw[flow] (lab.south) -- ++(0, -5.5) -| ([xshift=10pt]mg.north)
        node[pos=0.25, right=2pt, edgelbl] {Zapytania GUI};

    % Memgraph DB -> Wolumen trwały
    \draw[flow, dashed] (mg.east) -- (vol.west)
        node[midway, above=2pt, font=\scriptsize\bfseries] {Zapis};

\end{tikzpicture}%
}
\caption{Architektura środowiska kontenerowego Docker Compose. Schemat prezentuje separację sieciową na poziomie \texttt{frontend-net} oraz \texttt{backend-net}, rolę kontenera \texttt{backend} jako punktu styku oraz rozgraniczenie punktów komunikacji dla przeglądarki i mechanizmu Server-Side Rendering (SSR).}
\label{fig:docker_architektura_full}
\end{figure}

\newpage

\subsubsection{Obrazy kontenerowe i pliki Dockerfile}
% [X] 1. Analiza krok po kroku Dockerfile dla frontendu i backendu.
% [X] 2. Dobre praktyki: użytkownicy bez uprawnień root (node, appuser).
% [X] 3. Optymalizacja warstw (layer caching) przy instalacji zależności.
% [X] 4. Zmienne środowiskowe Pythona (wyłączenie .pyc, unbuffered logs). 

Kontenery \texttt{frontend} i \texttt{backend} uruchamiane są z obrazów
budowanych na podstawie plików \texttt{Dockerfile}, natomiast usługi
\texttt{memgraph-db} i \texttt{memgraph-lab} korzystają z gotowych, oficjalnych
obrazów. Instrukcje pliku \texttt{Dockerfile} wykonywane są sekwencyjnie podczas
budowania obrazu, a każda z nich tworzy kolejną warstwę. Oba obrazy są obecnie
skonfigurowane do pracy deweloperskiej (serwer Vite oraz Uvicorn z opcją
\texttt{--reload}), z zachowaniem zasad bezpieczeństwa (brak uprawnień konta
\texttt{root}) oraz optymalizacji pamięci podręcznej warstw
(ang. \emph{layer caching}):

\begin{itemize}
    \item frontend -- obraz składa się kolejno z:
    \begin{itemize}
        \item zdefiniowania obrazu bazowego \texttt{node:lts-slim},
        \item ustawienia katalogu roboczego \texttt{/app} i przekazania go na
          własność wbudowanemu użytkownikowi \texttt{node},
        \item skopiowania plików manifestu zależności (\texttt{package.json},
          \texttt{package-lock.json}) z nadaniem ich własności użytkownikowi
          \texttt{node} (\texttt{COPY --chown=node:node}),
        \item przełączenia na użytkownika \texttt{node}, dzięki czemu kolejne
          kroki i sama aplikacja nie działają z uprawnieniami konta \texttt{root},
        \item zainstalowania zależności menedżerem pakietów npm (oddzielenie
          tego kroku od kopiowania kodu źródłowego pozwala wykorzystać pamięć
          podręczną warstw -- zależności są instalowane ponownie tylko po zmianie
          plików manifestu),
        \item skopiowania kodu źródłowego, również z własnością użytkownika
          \texttt{node},
        \item ustawienia zmiennej \texttt{PORT} i udokumentowania portu
          \texttt{5173} (\texttt{EXPOSE}),
        \item zdefiniowania polecenia startowego uruchamiającego serwer
          deweloperski Vite (\texttt{npm run dev}) nasłuchujący na wszystkich
          interfejsach sieciowych.
    \end{itemize}
    \item backend -- obraz składa się kolejno z:
    \begin{itemize}
        \item zdefiniowania obrazu bazowego \texttt{python:3.12-slim},
        \item ustawienia zmiennych środowiskowych:
          \texttt{PYTHONDONTWRITEBYTECODE} (wyłączenie generowania plików
          \texttt{.pyc}), \texttt{PYTHONUNBUFFERED} (wypisywanie logów bez
          buforowania), \texttt{PIP\_NO\_CACHE\_DIR} (brak pamięci podręcznej
          menedżera pip, a więc mniejszy obraz) oraz
          \texttt{WATCHFILES\_FORCE\_POLLING} (wykrywanie zmian plików przez
          cykliczne odpytywanie, zapewniające niezawodne działanie opcji
          \texttt{--reload} w kontenerze),
        \item ustawienia katalogu roboczego \texttt{/app},
        \item utworzenia użytkownika \texttt{appuser},
        \item skopiowania pliku \texttt{requirements.txt} i zainstalowania
          zależności -- przed skopiowaniem kodu źródłowego, aby warstwa
          z zależnościami była ponownie wykorzystywana, dopóki plik wymagań się
          nie zmieni,
        \item skopiowania kodu źródłowego z nadaniem własności użytkownikowi
          \texttt{appuser} (\texttt{COPY --chown}),
        \item przełączenia na użytkownika \texttt{appuser} (zgodnie z zasadą
          najmniejszych uprawnień),
        \item udokumentowania portu \texttt{8000} (\texttt{EXPOSE}),
        \item zdefiniowania polecenia startowego uruchamiającego serwer Uvicorn
          z opcją \texttt{--reload}.
    \end{itemize}
\end{itemize}

\subsubsection{Mechanizmy integracji i konfiguracja środowiska deweloperskiego}

% [X] 1. Problem dwóch kontekstów sieciowych (SSR vs WebSocket):
%    - API_INTERNAL_URL (http://backend:8000 wewnątrz sieci Dockera)                                       %    - PUBLIC_API_BASE_URL (http://localhost:8000 dla przeglądarki)                                        % [X] 2. Automatyzacja pracy z Docker Compose Watch (develop.watch):
%    - synchronizacja plików w locie,                                                                      %    - automatyczny rebuild po zmianie package.json.                                                       % [X] 3. Kolejność uruchamiania i synchronizacja startu:
%    - ograniczenia dyrektywy depends_on (kontener wystartował != baza przyjmuje zapytania),               %    - mechanizm retry w backendzie przed seedowaniem grafu.

Oparcie architektury na kontenerach wymagało dostosowania konfiguracji
środowiska deweloperskiego oraz warstw serwerowej i interfejsowej. Konfiguracja
Docker Compose obejmuje mechanizmy zapewniające spójność komunikacji sieciowej,
automatyzację procesu wytwarzania oprogramowania oraz określoną kolejność
inicjalizacji usług.

\paragraph{Kontekst sieciowy}
Kluczowym wyzwaniem architektonicznym aplikacji frontendowej opartej na frameworku SvelteKit było złączenie w funkcjonalny sposób procesów wykonywanych po stronie serwera aplikacji oraz elementów renderujących po stronie przeglądarki klienta. Sytuacja ta generuje problem podwójnego kontekstu sieciowego:

\begin{itemize}                                                                                            \item \textbf{Kontekst wewnętrzny kontenera (SSR):} Podczas wykonywania funkcji ładujących na serwerze, proces Node.js działa wewnątrz izolowanej sieci mostkowej Dockera (\texttt{frontend-net}). W tym środowisku zapytania HTTP do backendu muszą być kierowane na adres wykorzystujący wewnętrzny resolver DNS Dockera, tj. \texttt{http://backend:8000}. Odwołanie się z wnętrza kontenera do adresu \texttt{localhost:8000} skutkowałoby błędem połączenia, gdyż adres pętli zwrotnej (ang. \emph{loopback}) odnosiłby się do samego kontenera frontendu, na którym usługa FastAPI nie nasłuchuje.\item \textbf{Kontekst zewnętrzny przeglądarki (CSR i WebSocket):} Kod wykonywany bezpośrednio w przeglądarce użytkownika (np. nawiązywanie stałego połączenia WebSocket dla podglądu symulacji pociągów w czasie rzeczywistym lub rezerwowe odpytywanie REST) działa na maszynie gospodarza. Przeglądarka nie posiada dostępu do wewnętrznej przestrzeni adresowej Dockera i nie rozpoznaje nazwy hosta \texttt{backend}. Wszelkie połączenia klienckie muszą być zatem kierowane na port wyeksponowany na maszynie gospodarza, czyli \texttt{http://localhost:8000} (lub odpowiednio \texttt{ws://localhost:8000/ws/live}). 
\end{itemize} Problem ten rozwiązano poprzez rozdzielenie konfiguracji na dwie zmienne środowiskowe: \texttt{API\_INTERNAL\_URL} (wykorzystywaną po stronie SSR i zabezpieczoną przed wyciekiem do klienta za pośrednictwem modułu \texttt{\$env/dynamic/private}) oraz \texttt{PUBLIC\_API\_BASE\_URL} (przekazywaną do klienta przez \texttt{\$env/dynamic/public}). Dzięki temu kod frontendu nigdy nie komunikuje się z bazą grafową bezpośrednio -- wszystkie zapytania przechodzą przez backend.  

\paragraph{Automatyzacja cyklu rozwoju oprogramowania za pomocą Docker Compose Watch} W projekcie wykorzystano mechanizm \texttt{develop.watch} wprowadzony w Docker Compose. Zastosowano dwutorową strategię synchronizacji:  

\begin{enumerate} 
    \item \textbf{Synchronizacja zmian w locie (\texttt{action: sync}):} Wszelkie modyfikacje kodu źródłowego w katalogach \texttt{./frontend} oraz \texttt{./backend} są natychmiast replikowane do ścieżki roboczej kontenera (\texttt{/app}). Z procesu synchronizacji wykluczono pliki generowane i specyficzne dla platformy (\texttt{node\_modules/}, \texttt{.svelte-kit/}, \texttt{.venv/}, \texttt{\_\_pycache\_\_}). Dzięki integracji z mechanizmem \textit{Hot Module Replacement} w serwerze Vite oraz flagą \texttt{--reload} w serwerze Uvicorn, zmiany w interfejsie i logice biznesowej są widoczne w czasie rzeczywistym bez konieczności restartu kontenerów. 
    \item \textbf{Automatyczna przebudowa obrazu (\texttt{action: rebuild}):}
    W przypadku modyfikacji pliku definiującego zależności frontendu
    (\texttt{package.json}) prosta synchronizacja plików jest niewystarczająca.
    Dyrektywa \texttt{rebuild} automatycznie inicjuje ponowne zbudowanie obrazu
    kontenera i zaktualizowanie środowiska uruchomieniowego po zapisaniu pliku,
    co eliminuje konieczność ręcznego wywoływania poleceń budowania przez
    programistę. W przypadku backendu reguła ta nie została skonfigurowana,
    dlatego zmiana pliku \texttt{requirements.txt} wymaga ręcznego przebudowania
    obrazu poleceniem \texttt{docker compose up --build backend}.                    
\end{enumerate}                       

\paragraph{Kolejność uruchamiania usług i synchronizacja startu}
Problemem było zapewnienie niezawodnej sekwencji uruchamiania powiązanych usług. W środowiskach kontenerowych typowym błędem jest zakładanie, że dyrektywa \texttt{depends\_on} gwarantuje gotowość usług zależnych. Standardowa dyrektywa \texttt{depends\_on} informuje jedynie demona Dockera o konieczności zainicjowania procesu kontenera bazy danych przed kontenerem backendu. Nie jest to równoważne z gotowością samej aplikacji do obsługi zapytań. Silnik bazy Memgraph po uruchomieniu kontenera wymaga od kilku do kilkunastu sekund na załadowanie modułów analitycznych MAGE, alokację struktur w pamięci RAM i otwarcie portu protokołu Bolt (\texttt{7687}). Próba natychmiastowego wykonania zapytań przez backend prowadziła do błędu \texttt{ServiceUnavailable} i awaryjnego zatrzymania kontenera. Problem wyeliminowano poprzez implementację mechanizmu aktywnego ponawiania prób w cyklu życia aplikacji FastAPI:

\begin{itemize}
    \item Funkcja \texttt{connect\_with\_retry} w pętli tworzy połączenie
    i sprawdza jego dostępność; w razie wyjątku \texttt{ServiceUnavailable}
    ponawia próbę co 2~s, aż baza będzie gotowa.
    \item Po nawiązaniu połączenia moduł startowy sprawdza liczbę węzłów
    \texttt{:Station}. Jeżeli baza jest pusta, uruchamiany jest skrypt
    zasilający (\emph{seed}), który tworzy sieć kolejową oraz flotę pociągów;
    w przeciwnym razie krok ten jest pomijany, co zapobiega nadpisaniu danych
    przy kolejnych restartach.
    \item Następnie przy każdym starcie zakładane są brakujące elementy:
    startowe scenariusze rozkładów, ograniczenia unikalności (adresy e-mail,
    nazwy ról), role systemowe oraz konto administratora -- wyłącznie wtedy,
    gdy jeszcze nie istnieją.
\end{itemize}
Dzięki temu środowisko startuje w sposób w pełni autonomiczny, odporny na asynchroniczne opóźnienia rozruchu poszczególnych komponentów.   

\subsection{Warstwa prezentacyjna}
% Backend oraz to jak cały projekt komunikuję się z sobą
\subsection{Warstwa komunikacyjno-zarządzająca}
Warstwa komunikacyjno-zarządzająca (ang. \emph{backend}) stanowi najważniejsze
ogniwo całej symulacji -- zachodzą w niej procesy decyzyjne i działa w niej
silnik symulacji. Typowe rozwiązania backendowe opierają się głównie na
operacjach CRUD, natomiast opisywany komponent łączy cechy bezstanowego
interfejsu API, asynchronicznego serwera rozgłaszającego stan symulacji oraz
autonomicznego silnika symulacji. Do zadań tej warstwy należą:
\begin{itemize}
    \item koordynacja w czasie rzeczywistym ruchu floty pociągów na podstawie rzeczywistych parametrów infrastruktury kolejowej węzła śląskiego,
    \item wyznaczanie optymalnych czasowo tras przejazdu oraz automatyczne planowanie objazdów w warunkach dynamicznie zmieniającego się stanu sieci (algorytm $A^*$),
    \item modelowanie losowych incydentów infrastrukturalnych i taborowych z wykorzystaniem jednorodnego procesu Poissona,
    \item dwutorowa wymiana danych z klientami: transakcyjna komunikacja sterująca poprzez protokół REST (HTTP) oraz rozgłaszanie stanu symulacji przez strumień WebSocket,
    \item uwierzytelnianie i autoryzacja użytkowników w oparciu o model kontroli
    dostępu opartej na rolach (\emph{Role-Based Access Control}, RBAC), w którym
    użytkownicy (\texttt{:User}) oraz role wraz z listami uprawnień
    (\texttt{:Role}) przechowywane są jako węzły bazy grafowej.
\end{itemize}

Czytelność oraz łatwość dodawania nowych rozwiązań były nadrzędnymi celami
podczas budowania architektury backendu. Układ katalogów i plików
implementacyjnych przedstawiono na rysunku~\ref{fig:backend_tree}.

\begin{figure}[htbp]
\centering
\begin{minipage}{\textwidth}
\setlength{\parindent}{0pt}
\singlespacing
\footnotesize
\DTsetlength{0.2em}{1.2em}{0.2em}{0.4pt}{1.6pt}
\setlength{\DTbaselineskip}{11pt}
\dirtree{%
.1 backend/.
.2 Dockerfile\DTcomment{budowanie obrazu Python 3.12-slim}.
.2 requirements.txt\DTcomment{zależności Pythona (FastAPI, neo4j, uvicorn\ldots)}.
.2 pytest.ini\DTcomment{konfiguracja środowiska testowego}.
.2 README.md\DTcomment{podstawowa dokumentacja backendu}.
.2 app/.
.3 main.py\DTcomment{punkt wejścia FastAPI, CORS, routery}.
.3 core/\DTcomment{moduły jądra systemu}.
.4 config.py\DTcomment{parametry symulacji, wagi awarii}.
.4 lifespan.py\DTcomment{cykl życia: start/stop bazy i symulacji}.
.4 ws\_manager.py\DTcomment{pula połączeń WebSocket i rozgłaszanie}.
.4 haversine.py\DTcomment{wzór Haversine'a, dopuszczalna heurystyka A*}.
.3 api/\DTcomment{warstwa transportowa HTTP i WebSocket}.
.4 dependencies.py\DTcomment{wstrzykiwanie zależności, RBAC}.
.4 routes/\DTcomment{definicje końcówek API (cienkie kontrolery)}.
.5 admin.py\DTcomment{zarządzanie użytkownikami i rolami}.
.5 auth.py\DTcomment{logowanie, rejestracja, sesja gościa, profil /me}.
.5 events.py\DTcomment{pobieranie zdarzeń, ręczne zgłaszanie incydentów}.
.5 health.py\DTcomment{prosty healthcheck (/api/hello)}.
.5 live.py\DTcomment{główny kanał WebSocket (/ws/live)}.
.5 network.py\DTcomment{graf sieci kolejowej: stacje i odcinki}.
.5 routing.py\DTcomment{wyznaczanie najszybszej trasy A*}.
.5 scenarios.py\DTcomment{CRUD scenariuszy i uruchamianie rozkładów}.
.5 simulation.py\DTcomment{tempo, pauza, reset floty}.
.5 trains.py\DTcomment{migawka stanu pociągów (rezerwowy REST)}.
.3 schemas/\DTcomment{kontrakty danych i walidacja Pydantic v2}.
.4 auth.py\DTcomment{DTO autoryzacji, użytkowników i ról}.
.4 event.py\DTcomment{DTO zdarzeń losowych i incydentów}.
.4 network.py\DTcomment{DTO stacji, odcinków i kierunków}.
.4 routing.py\DTcomment{DTO wyników wyznaczania trasy}.
.4 scenario.py\DTcomment{DTO planu pociągu i scenariuszy rozkładu}.
.4 simulation.py\DTcomment{DTO nastaw tempa symulacji}.
.4 train.py\DTcomment{DTO pociągów i migawki floty}.
.3 services/\DTcomment{logika biznesowa i algorytmy}.
.4 auth\_service.py\DTcomment{haszowanie PBKDF2, tokeny HMAC, RBAC}.
.4 event\_service.py\DTcomment{awarie (proces Poissona)}.
.4 network\_service.py\DTcomment{odczyt grafu, agregacja odcinków}.
.4 routing\_service.py\DTcomment{A*: stan odcinków, typ pociągu}.
.4 scenario\_service.py\DTcomment{rozkłady, falowe wjazdy, pauza}.
.4 simulation\_engine.py\DTcomment{asynchroniczna pętla symulacji}.
.4 train\_service.py\DTcomment{ruch, postoje, wykolejenia, rerouting}.
.2 db/\DTcomment{inicjalizacja i dane początkowe}.
.3 seed.py\DTcomment{skrypt zasilający bazę węzłami i relacjami}.
.3 README.md\DTcomment{dokumentacja struktury bazy grafowej}.
.2 tests/\DTcomment{testy automatyczne pytest (9 modułów)}.
}
\end{minipage}
\caption{Struktura katalogów warstwy backendu}
\label{fig:backend_tree}
\end{figure}

\subsubsection{Architektura logiczna, wzorce projektowe i organizacja kodu}
Warstwa backendu została zrealizowana w oparciu o architekturę warstwową
(ang. \emph{Layered Architecture}), zgodnie z zasadą separacji odpowiedzialności
(ang. \emph{Separation of Concerns}). Kod podzielono tak, aby nie powstawały
zależności cykliczne, a testy jednostkowe były łatwe do przeprowadzenia.

    \paragraph{Separacja odpowiedzialności (Separation of Concerns)}
    Moduły zlokalizowane w \texttt{app/api/routes} pełnią rolę cienkich
    kontrolerów. Ich zadaniem jest przyjmowanie żądań HTTP lub ramek WebSocket,
    walidacja parametrów wejściowych, wywołanie odpowiedniej metody z warstwy
    logiki biznesowej (\texttt{app/services/}) oraz zwrócenie odpowiedzi
    z właściwym kodem statusu. W kontrolerach nie implementuje się zapytań do
    bazy danych ani algorytmów -- cała logika obliczeniowa i zarządzająca
    znajduje się w warstwie serwisowej.
    \paragraph{Wzorzec wstrzykiwania zależności (Dependency Injection)}
    \paragraph{Kontrakty danych i deterministyczna walidacja (Pydantic v2)}
\subsubsection{Dwutorowy interfejs komunikacyjny (REST API oraz WebSockets)}
    \paragraph{Sterująco-transakcyjne API REST}
    \paragraph{Kanał telemetryczny WebSocket i zarządca połączeń}
    \paragraph{Separacja kontekstów i zabezpieczenie nagłówka Origin}
\subsubsection{Autonomiczny silnik symulacji i model współbieżności}
    \paragraph{Cykl życia aplikacji i asynchroniczna pętla symulacji}
    \paragraph{Izolacja wątkowa (asyncio.to\_thread) i synchronizacja stanowa}
    \paragraph{Chronologia pojedynczego kroku symulacji (Tick Lifecycle)}
    \paragraph{Model ruchu pociągów, skala czasu i obsługa pauzy}
\subsubsection{Algorytm wyznaczania tras (A*) i dynamiczny rerouting}
    \paragraph{Model grafu i funkcja kosztu rzeczywistego}
    \paragraph{Dopuszczalna heurystyka ortodromiczna (wzór Haversine'a)}
    \paragraph{Świadomość kierunkowa oraz filtracja stanowa krawędzi}
    \paragraph{Dynamiczne przeliczanie tras alternatywnych (rerouting w locie)}
\subsubsection{Stochastyczny model zakłóceń i obsługa sytuacji kryzysowych}
    \paragraph{Generowanie zdarzeń losowych w oparciu o proces Poissona}
    \paragraph{Typologia incydentów i mechanika propagacji skutków}
    \paragraph{Deterministyczna rezolucja awarii i manualna ingerencja dyspozytorska}
\subsubsection{Bezpieczeństwo, kontrola dostępu i model RBAC}
    \paragraph{Kryptograficzna ochrona danych uwierzytelniających (PBKDF2)}
    \paragraph{Bezstanowe tokeny sesyjne z podpisem HMAC-SHA256}
    \paragraph{Model RBAC i reguły ochrony przed samoblokadą}
\subsubsection{Strategia testowania i weryfikacja poprawności systemu}
    \paragraph{Organizacja zestawu testowego (pytest)}
    \paragraph{Weryfikacja kluczowych scenariuszy brzegowych i algorytmicznych}

% Opisanie architektury bazy danych
% Źródła czemu ten rodzaj bazy lepszy, czemu ją wybraliśmy, jak powstała baza, źródła skąd powstawała, 
\subsection{Warstwa bazodanowa}
\section{Stos technologiczny}
W niniejszym rozdziale szczegółowo opisano narzędzia oraz technologie wybrane do
realizacji projektu interaktywnego symulatora ruchu kolejowego. Dobór stosu
technologicznego był podyktowany potrzebą zapewnienia wysokiej wydajności obliczeniowej
systemu, szczególnie w zakresie przetwarzania danych grafowych, oraz koniecznością
stworzenia responsywnego i intuicyjnego interfejsu użytkownika. Architektura systemu
została podzielona na niezależne warstwy: kliencką, serwerową oraz bazodanową, co
pozwala na zachowanie czystości kodu i łatwą rozbudowę poszczególnych modułów
w przyszłości.

\subsection{Warstwa kliencka}

Warstwa kliencka (ang. \emph{frontend}) odpowiada za wizualizację siatki połączeń
kolejowych na schematycznej mapie odwzorowującej położenie stacji oraz prezentację stanów pociągów w czasie
rzeczywistym. Kluczowym wyzwaniem w tej warstwie było zapewnienie płynności animacji
przy dużej liczbie przemieszczających się obiektów, przy jednoczesnym utrzymaniu
niskiego zapotrzebowania na zasoby sprzętowe użytkownika.

\subsubsection{Svelte i SvelteKit}

Głównym narzędziem wybranym do budowy interfejsu użytkownika jest kompilator Svelte.
W odróżnieniu od tradycyjnych bibliotek, Svelte przenosi większość pracy na etap
budowania aplikacji, generując wysoce zoptymalizowany kod w czystym języku
JavaScript~\cite{svelte}. Pozwala to na uniknięcie narzutu obliczeniowego związanego
z ciągłym przeliczaniem wirtualnego modelu obiektu dokumentu (ang. \emph{Virtual
Document Object Model -- VDOM}). W przypadku symulatora, w którym pozycje wielu
pociągów muszą być stale aktualizowane z ułamkową precyzją, brak tego narzutu
drastycznie poprawia płynność działania aplikacji. Jako nadrzędny szkielet
aplikacyjny (ang. \emph{framework}) zastosowano SvelteKit, co ułatwiło zarządzanie
architekturą projektu oraz trasowaniem (ang. \emph{routing}).

\subsubsection{Wizualizacja sieci w formacie SVG}

Do wizualizacji sieci kolejowej wykorzystano skalowalną grafikę wektorową
(ang. \emph{Scalable Vector Graphics} -- SVG)~\cite{svg}, renderowaną
bezpośrednio przez komponent Svelte (\texttt{NetworkGraph.svelte}), bez
zewnętrznych bibliotek mapowych. Położenie stacji wyznaczane jest przez liniowe
przeskalowanie współrzędnych geograficznych do wymiarów obszaru rysowania, co
zachowuje wzajemny układ stacji. Pozycja pociągu jest interpolowana liniowo
między stacjami na podstawie wskaźnika postępu przesyłanego z serwera.
Przesuwanie i przybliżanie widoku zrealizowano samodzielnie z użyciem zdarzeń
wskaźnika (ang. \emph{pointer events}) i transformacji współrzędnych SVG.

Rezygnacja z mapy podkładowej była świadomą decyzją: schematyczny widok jest
czytelniejszy przy dużej liczbie znaczników, nie wymaga pobierania kafelków
z zewnętrznych serwerów i działa bez dostępu do internetu, co jest spójne
z założeniem autonomicznej piaskownicy. Każdy element sieci (stacja, odcinek,
pociąg) jest elementem drzewa DOM, dzięki czemu wybór i podświetlanie elementów
nie wymagają dodatkowej warstwy logiki.

\subsubsection{TypeScript}

Całość logiki po stronie klienta została napisana w języku TypeScript, który wprowadza
do JavaScriptu silne typowanie statyczne~\cite{typescript}. Ścisła integracja
TypeScriptu ze środowiskiem Svelte pozwala na rygorystyczne definiowanie struktur
danych przesyłanych z serwera (np. węzłów grafu, parametrów fizycznych składów).
Wczesne wykrywanie błędów na etapie kompilacji znacząco zwiększa stabilność kodu
i ułatwia zarządzanie rosnącą złożonością systemu informatycznego.

\subsubsection{Style interfejsu i Tailwind CSS}

Wygląd interfejsu zdefiniowano głównie w stylach lokalnych komponentów Svelte
(bloki \texttt{<style>}), których zasięg jest automatycznie ograniczany do
danego komponentu, dzięki czemu reguły jednego panelu nie wpływają na pozostałe.
W projekcie skonfigurowano również bibliotekę Tailwind CSS~\cite{tailwind},
opartą na klasach użytkowych (ang. \emph{utility-first}), wykorzystywaną do
stylów globalnych oraz pojedynczych elementów.

\subsubsection{Licencje}

Większość wykorzystanych technologii jest udostępniana na licencjach o otwartym
kodzie źródłowym (ang. \emph{open-source}), w tym na permisywnych licencjach MIT
(Svelte, SvelteKit, Tailwind CSS, FastAPI) oraz Apache~2.0 (TypeScript,
sterownik \texttt{neo4j}). Wyjątki stanowią biblioteka animacji GSAP,
udostępniana bezpłatnie na własnej licencji producenta~\cite{gsaplicense}, oraz
baza Memgraph w wersji Community, objęta licencją Business Source
License~1.1~\cite{memgraphlicense}, która pozwala na bezpłatne użycie,
z ograniczeniami dotyczącymi komercyjnego udostępniania bazy jako usługi.
Umożliwia to legalne wykorzystanie całego stosu w projekcie akademickim.

\subsection{Warstwa serwerowa}

Warstwa serwerowa (ang. \emph{backend}) stanowi główne centrum obliczeniowe i decyzyjne
symulatora, realizujące zaawansowaną logikę biznesową, w tym wyznaczanie tras oraz
symulację ruchu pociągów. Warstwa ta musi jednocześnie wykonywać cykliczne kroki symulacji, obsługiwać
żądania użytkowników i rozsyłać stan symulacji do wszystkich podłączonych
przeglądarek. Z tego względu oparto ją na modelu asynchronicznym,
zoptymalizowanym pod kątem operacji wejścia-wyjścia.

\subsubsection{Python i FastAPI}

Jako główny język programowania po stronie serwera wybrano Pythona, co było podyktowane
jego wszechstronnością oraz pozycją rynkowego standardu w środowiskach inżynierskich.
Czytelna składnia i bogata biblioteka standardowa (m.in. moduły \texttt{heapq},
\texttt{math} i \texttt{asyncio}) pozwoliły zaimplementować algorytm A* oraz
silnik symulacji samodzielnie, bez zewnętrznych bibliotek grafowych, co daje
pełną kontrolę nad funkcją kosztu i heurystyką. W roli nadrzędnego szkieletu programistycznego (ang. \emph{framework})
wykorzystano FastAPI, czyli nowoczesny, wydajny framework oparty na standardzie ASGI. Narzędzie to natywnie obsługuje asynchroniczność oraz
automatycznie generuje interaktywną dokumentację interfejsu programowania aplikacji
(ang. \emph{Application Programming Interface -- API}) w standardzie OpenAPI, co znacząco
usprawniło testowanie komunikacji pomiędzy serwerem a aplikacją kliencką~\cite{fastapi}.

\subsubsection{Warstwa przetwarzania asynchronicznego}

Pozycje pociągów muszą być aktualizowane w stałych odstępach czasu, a serwer
musi przy tym pozostać responsywny. System wykorzystuje wbudowaną w język Python
pętlę zdarzeń (ang. \emph{event loop})~\cite{pythonasyncio}, w której działa
zadanie symulacji uruchamiane co sekundę. Ponieważ sterownik bazy danych działa
synchronicznie, obliczenia pojedynczego kroku przenoszone są do osobnego wątku
(\texttt{asyncio.to\_thread}), dzięki czemu nie blokują pętli zdarzeń.
Aplikacja może zatem jednocześnie obsługiwać żądania użytkowników (np.
zgłoszenie awarii) i bez przerw przesyłać zaktualizowane pozycje pociągów do
interfejsu.

\subsection{Warstwa bazodanowa}

Warstwa trwałości danych odpowiada za przechowywanie oraz błyskawiczne udostępnianie
kluczowych informacji o strukturze sieci kolejowej. Z uwagi na fakt, że mapa połączeń
transportowych jest w swojej ostatecznej formie złożonym grafem, tradycyjne relacyjne
bazy danych są mniej odpowiednie ze względu na wysoki koszt obliczeniowy zapytań
o wielokrotnie zagnieżdżone relacje~\cite{graphdatabases}.

\subsubsection{Memgraph}

Fundamentem systemu zarządzania danymi w projekcie jest grafowa baza danych
Memgraph~\cite{memgraph}. Przechowuje ona informacje w postaci wierzchołków
reprezentujących stacje kolejowe oraz krawędzi odzwierciedlających łączące je szlaki,
co stanowi bardzo naturalne i intuicyjne odwzorowanie fizycznej infrastruktury torowej.
Kluczową przewagą bazy Memgraph jest jej architektura oparta na przetwarzaniu danych
w całości w pamięci operacyjnej komputera (ang. \emph{in-memory}). Cała sieć (57 stacji i 142 skierowane relacje \texttt{TRACK}) jest odczytywana
z bazy przy każdym wyznaczaniu trasy, a samo przeszukiwanie algorytmem A*
odbywa się w warstwie aplikacji. Memgraph obsługuje język zapytań
openCypher~\cite{cypher}, w którym warunki takie jak pominięcie zamkniętych
odcinków wyraża się bezpośrednio w zapytaniu, np.
\texttt{WHERE r.status IN ['active', 'restricted']}.

\subsection{Konteneryzacja}

W celu zapewnienia wysokiej powtarzalności środowiska uruchomieniowego oraz ułatwienia
procesu wdrożenia (ang. \emph{deployment}) aplikacji na różnych infrastrukturach
serwerowych, w projekcie zdecydowano się na zastosowanie technologii konteneryzacji.
Podejście to pozwala na hermetyczne odizolowanie poszczególnych warstw systemu --
klienckiej, serwerowej oraz bazodanowej -- od bezpośredniej konfiguracji systemu
operacyjnego maszyny hosta~\cite{docker}. Dzięki takiemu rozwiązaniu wyeliminowano
powszechny w inżynierii oprogramowania problem niezgodności wersji bibliotek systemowych
oraz ułatwiono skalowanie całego rozwiązania. Konteneryzacja sprzyja również utrzymaniu
czystości w środowisku deweloperskim, gdyż wszystkie niezbędne zależności są instalowane
wyłącznie wewnątrz odizolowanych jednostek, co gwarantuje spójność działania aplikacji
niezależnie od platformy, na której zostanie uruchomiona.

\subsubsection{Docker}

Podstawowym narzędziem wykorzystanym do wirtualizacji na poziomie systemu operacyjnego
jest Docker. Pozwala on na zamknięcie każdej części aplikacji w niezależnym, lekkim
środowisku uruchomieniowym zwanym kontenerem (ang. \emph{container}). W kontekście symulatora kolejowego przygotowano własne obrazy
(ang. \emph{images}) dla aplikacji serwerowej napisanej w języku Python oraz
warstwy klienckiej opartej na technologii Svelte; każdy z nich zawiera komplet
bibliotek i narzędzi niezbędnych do pracy danego modułu. Baza danych Memgraph
i narzędzie Memgraph Lab uruchamiane są z gotowych, oficjalnych obrazów. Wykorzystanie Dockera pozwala
na szybkie odtworzenie całego stosu technologicznego na dowolnym urządzeniu, co jest
kluczowe podczas pracy zespołowej oraz przy późniejszych testach integracyjnych systemu.
Dodatkowo, mechanizm warstwowości obrazów Dockerowych optymalizuje czas budowania
aplikacji, co znacząco przyspiesza proces iteracyjnego wprowadzania poprawek w logice symulacji ruchu czy algorytmach dynamicznego wyznaczania objazdów.

\subsubsection{Docker Compose}

Do zarządzania architekturą składającą się z wielu współpracujących ze sobą kontenerów
wykorzystano narzędzie Docker Compose. Umożliwia ono zdefiniowanie całej infrastruktury
projektowej w postaci pojedynczego, deklaratywnego pliku konfiguracyjnego, co stanowi
realizację nowoczesnego paradygmatu infrastruktury jako kodu (ang. \emph{Infrastructure
as Code -- IaC})~\cite{infrastructureascode}. W pliku tym precyzyjnie określono zależności pomiędzy poszczególnymi usługami,
mapowanie portów sieciowych, dwie izolowane sieci wewnętrzne oraz nazwany
wolumin danych (ang. \emph{volume}), który zapewnia trwałość informacji
przechowywanych w bazie Memgraph nawet po usunięciu kontenera. Podział na sieci
sprawia, że interfejs użytkownika nie komunikuje się z bazą danych
bezpośrednio, lecz wyłącznie za pośrednictwem aplikacji serwerowej.
Uruchomienie całego systemu symulacyjnego sprowadza się do wykonania jednego
polecenia (\texttt{docker compose up}), co znacząco obniża próg wejścia dla
nowych osób pracujących nad projektem i minimalizuje ryzyko błędów
konfiguracyjnych.

\dodatkowo{Podsumowanie}

% TODO: 2-3 akapity: co zrealizowano, czy cele z rozdziału 2 zostały osiągnięte,
% najważniejsze wyniki (57 stacji, 71 odcinków, 52 pociągi, A*, zdarzenia losowe).

\subsection*{Kierunki dalszego rozwoju}
\addcontentsline{toc}{subsection}{Kierunki dalszego rozwoju}

Zrealizowany system stanowi podstawę do dalszej rozbudowy. Najważniejsze
kierunki rozwoju obejmują:

\begin{itemize}
\item blokadę zajętości odcinków (rezerwację odcinka dla jednego pociągu,
  odpowiednik samoczynnej blokady liniowej) wraz z rozstrzyganiem pierwszeństwa
  na odcinkach jednotorowych,
\item priorytety kategorii pociągów, np. ustępowanie pociągu towarowego
  pociągowi dalekobieżnemu na stacji z torem dodatkowym,
\item model dynamiki z przyspieszaniem i hamowaniem, wykorzystujący zapisane
  już parametry składów (\texttt{accel}, \texttt{decel}, masa),
\item model rozkładu jazdy z propagacją opóźnień i wskaźnikiem punktualności,
\item mapę podkładową, np. Leaflet~\cite{leaflet} z kafelkami
  OpenStreetMap~\cite{openstreetmap}.
\end{itemize}

%%%%%%%%%%%%%%%%%%%%%%%%%%% bibliografia -- recznie, wg Wzoru A (formularz Z4-PU12)
%% UWAGA: daty dostepu do stron WWW oraz dane pozycji nalezy zweryfikowac.
\begin{thebibliography}{30}

\bibitem{leaflet} Agafonkin V., Leaflet -- a JavaScript library for interactive maps, \url{https://leafletjs.com/} (dostęp: 30.06.2026).

\bibitem{agilemanifesto} Beck K. i in., Manifesto for Agile Software Development, 2001, \url{https://agilemanifesto.org/} (dostęp: 30.06.2026).

\bibitem{srk} Dąbrowa-Bajon M., Podstawy sterowania ruchem kolejowym. Funkcje, wymagania, zarys techniki, Oficyna Wydawnicza Politechniki Warszawskiej, Warszawa 2014.

\bibitem{docker} Docker Inc., Docker Documentation, \url{https://docs.docker.com/} (dostęp: 30.06.2026).

\bibitem{digitaltwin} Errandonea I., Beltrán S., Arrizabalaga S., Digital Twin for maintenance: A literature review, Computers in Industry, vol. 123, 2020, str. 103316.

\bibitem{gsaplicense} GreenSock, GSAP Standard License, \url{https://gsap.com/community/standard-license/} (dostęp: 28.09.2026).

\bibitem{astar1968} Hart P.E., Nilsson N.J., Raphael B., A Formal Basis for the Heuristic Determination of Minimum Cost Paths, IEEE Transactions on Systems Science and Cybernetics, vol. 4, nr 2, 1968, str. 100--107.

\bibitem{gtfsks} Koleje Śląskie sp. z o.o., Rozkład jazdy w formacie GTFS, wersja 2026.166 (kopia lustrzana: repozytorium gtfs-proxies/24-Koleje-Slaskie), 2026.

\bibitem{martin2017clean} Martin R.C., Clean Architecture: A Craftsman's Guide to Software Structure and Design, Prentice Hall, Boston 2017.

\bibitem{memgraphlicense} Memgraph, Business Source License 1.1, \url{https://github.com/memgraph/memgraph/blob/master/licenses/BSL.txt} (dostęp: 28.09.2026).

\bibitem{memgraph} Memgraph, Memgraph Documentation, \url{https://memgraph.com/docs} (dostęp: 30.06.2026).

\bibitem{kontener.microsoft} Microsoft, Co to jest kontener?, \url{https://azure.microsoft.com/pl-pl/resources/cloud-computing-dictionary/what-is-a-container} (dostęp: 11.09.2026).

\bibitem{typescript} Microsoft, TypeScript Documentation, \url{https://www.typescriptlang.org/docs/} (dostęp: 30.06.2026).

\bibitem{railwaymarket} Mordor Intelligence, Railway Management System Market -- Growth, Trends and Forecasts, Mordor Intelligence, 2024, \url{https://www.mordorintelligence.com/} (dostęp: 30.06.2026).

\bibitem{infrastructureascode} Morris K., Infrastructure as Code: Managing Servers in the Cloud, wyd. 2, O'Reilly Media, Sebastopol 2020.

\bibitem{cypher} Neo4j, Cypher Query Language Reference, \url{https://neo4j.com/docs/cypher-manual/} (dostęp: 30.06.2026).

\bibitem{openstreetmap} OpenStreetMap contributors, OpenStreetMap, \url{https://www.openstreetmap.org/} (dostęp: 30.06.2026).

\bibitem{plkregulamin} PKP Polskie Linie Kolejowe S.A., Regulamin sieci 2025/2026, załączniki 1, 2.1, 2.4, 2.6 i 2.18, \url{https://www.plk-sa.pl/} (dostęp: 30.08.2026).

\bibitem{pythonasyncio} Python Software Foundation, asyncio -- Asynchronous I/O. Python Documentation, \url{https://docs.python.org/3/library/asyncio.html} (dostęp: 30.06.2026).

\bibitem{fastapi} Ramírez S., FastAPI Documentation, \url{https://fastapi.tiangolo.com/} (dostęp: 30.06.2026).

\bibitem{graphdatabases} Robinson I., Webber J., Eifrem E., Graph Databases: New Opportunities for Connected Data, wyd. 2, O'Reilly Media, Sebastopol 2015.

\bibitem{haversine} Sinnott R.W., Virtues of the Haversine, Sky and Telescope, vol. 68, nr 2, 1984, str. 159.

\bibitem{sommerville2016} Sommerville I., Software Engineering, wyd. 10, Pearson, Boston 2016.

\bibitem{svelte} Svelte, Svelte Documentation, \url{https://svelte.dev/docs} (dostęp: 30.06.2026).

\bibitem{tailwind} Tailwind Labs, Tailwind CSS Documentation, \url{https://tailwindcss.com/docs} (dostęp: 30.06.2026).

\bibitem{svg} W3C, Scalable Vector Graphics (SVG) 2, \url{https://www.w3.org/TR/SVG2/} (dostęp: 28.09.2026).

\end{thebibliography}

\end{document}
